% from ega1-3-fr.tex
\oldpage[I]{106}
Il suffit d'appliquer (\hyperref[I.3.2.4-fr]{3.2.4}) à l'injection canonique
$S'\to S$.

\begin{theorem}[3.2.6]
\label{I.3.2.6-fr}
Étant donnés deux $S$-préschémas $X$, $Y$, il existe un produit
$X\times_S Y$.
\end{theorem}

Nous procéderons en plusieurs étapes.

\begin{lemma}[3.2.6.1]
\label{I.3.2.6.1-fr}
Soient $(Z,p,q)$ un produit de $X$ et $Y$, $U$, $V$ des parties ouvertes de
$X$, $Y$ respectivement. Si on pose
$W=p^{-1}(U)\cap q^{-1}(V)$, le triplet formé de $W$ et des restrictions de
$p$ et $q$ à $W$ \textup{(considérées comme des morphismes $W\to U$,
$W\to V$ respectivement)} est un produit de $U$ et $V$.
\end{lemma}

En effet, si $T$ est un $S$-préschéma, on peut identifier les $S$-morphismes
$T\to W$ et les $S$-morphismes $T\to Z$ appliquant $T$ dans $W$. Si alors
$g:T\to U$, $h:T\to V$ sont deux $S$-morphismes quelconques, on peut les
considérer comme des $S$-morphismes de $T$ dans $X$ et $Y$ respectivement, et
par hypothèse il y a donc un $S$-morphisme et un seul $f:T\to Z$ tel que
$g=p\circ f$, $h=q\circ f$. Comme $p(f(T))\subset U$, $q(f(T))\subset V$,
on a
\[
  f(T)\subset p^{-1}(U)\cap q^{-1}(V)=W,
\]
d'où notre assertion.

\begin{lemma}[3.2.6.2]
\label{I.3.2.6.2-fr}
Soient $Z$ un $S$-préschéma, $p:Z\to X$, $q:Z\to Y$ deux $S$-morphismes,
$(U_\alpha)$ un recouvrement ouvert de $X$, $(V_\lambda)$ un recouvrement
ouvert de $Y$. On suppose que pour tout couple $(\alpha,\lambda)$, le
$S$-préschéma
$W_{\alpha\lambda}=p^{-1}(U_\alpha)\cap q^{-1}(V_\lambda)$ et les
restrictions de $p$ et $q$ à $W_{\alpha\lambda}$ constituent un produit de
$U_\alpha$ et de $V_\lambda$. Alors $(Z,p,q)$ est un produit de $X$ et $Y$.
\end{lemma}

Montrons d'abord que, si $f_1$, $f_2$ sont deux $S$-morphismes $T\to Z$, les
relations $p\circ f_1=p\circ f_2$ et $q\circ f_1=q\circ f_2$ entraînent
$f_1=f_2$. En effet, $Z$ est réunion des $W_{\alpha\lambda}$, donc les
$f_1^{-1}(W_{\alpha\lambda})$ forment un recouvrement ouvert de $T$, et il
en est de même des $f_2^{-1}(W_{\alpha\lambda})$. En outre, on a
\[
\begin{aligned}
  f_1^{-1}(W_{\alpha\lambda})
  &=f_1^{-1}(p^{-1}(U_\alpha))\cap
    f_1^{-1}(q^{-1}(V_\lambda)) \\
  &=f_2^{-1}(p^{-1}(U_\alpha))\cap
    f_2^{-1}(q^{-1}(V_\lambda))
   =f_2^{-1}(W_{\alpha\lambda})
\end{aligned}
\]
par hypothèse, et tout revient à voir que les restrictions de $f_1$ et $f_2$
à $f_1^{-1}(W_{\alpha\lambda})=f_2^{-1}(W_{\alpha\lambda})$ sont
identiques pour tout couple d'indices. Mais comme ces restrictions peuvent
être considérées comme des $S$-morphismes de
$f_1^{-1}(W_{\alpha\lambda})$ dans $W_{\alpha\lambda}$, notre assertion
résulte des hypothèses et de la déf. (\hyperref[I.3.2.1-fr]{3.2.1}).

Supposons maintenant donnés deux $S$-morphismes $g:T\to X$, $h:T\to Y$.
Posons
$T_{\alpha\lambda}=g^{-1}(U_\alpha)\cap h^{-1}(V_\lambda)$~; les
$T_{\alpha\lambda}$ forment un recouvrement ouvert de $T$. Par hypothèse,
il existe un $S$-morphisme $f_{\alpha\lambda}$ tel que
$p\circ f_{\alpha\lambda}$ et $q\circ f_{\alpha\lambda}$ soient les
restrictions respectives de $g$ et $h$ à $T_{\alpha\lambda}$. En outre,
montrons que les restrictions de $f_{\alpha\lambda}$ et
$f_{\beta\mu}$ à
$T_{\alpha\lambda}\cap T_{\beta\mu}$ coïncident, ce qui achèvera de
prouver (\hyperref[I.3.2.6.2-fr]{3.2.6.2}). Or, les images de
$T_{\alpha\lambda}\cap T_{\beta\mu}$ par $f_{\alpha\lambda}$ et par
$f_{\beta\mu}$ sont contenues dans
$W_{\alpha\lambda}\cap W_{\beta\mu}$ par définition. Comme
\[
  W_{\alpha\lambda}\cap W_{\beta\mu}
  =p^{-1}(U_\alpha\cap U_\beta)\cap
   q^{-1}(V_\lambda\cap V_\mu),
\]
il résulte de (\hyperref[I.3.2.6.1-fr]{3.2.6.1}) que
$W_{\alpha\lambda}\cap W_{\beta\mu}$ et les restrictions à ce préschéma
de $p$ et $q$ constituent un produit de
$U_\alpha\cap U_\beta$ et de $V_\lambda\cap V_\mu$. Comme
$p\circ f_{\alpha\lambda}$ et $p\circ f_{\beta\mu}$ coïncident dans
$T_{\alpha\lambda}\cap T_{\beta\mu}$ et qu'il en est de même de
$q\circ f_{\alpha\lambda}$ et $q\circ f_{\beta\mu}$, on voit que
$f_{\alpha\lambda}$ et $f_{\beta\mu}$ coïncident dans
$T_{\alpha\lambda}\cap T_{\beta\mu}$, c.q.f.d.

